\naslov{Teorija upodobitev}

\begin{definicija}
  \pojem{$R$-algebra} je levi $R$-modul $A$, ki je hkrati kolobar z enoto, in
  velja $r(a_1 a_2) = (r a_1) a_2 = a_1 (r a_2)$.
\end{definicija}

\begin{primer}
  Če je $R$ komutativen kolobar z enico in $G$ grupa, je prosti $R$-modul nad
  $G$, ki ga označimo z $RG$, $R$-algebra, ki ji pravimo \pojem{grupna algebra}.
\end{primer}

\begin{definicija}
  Naj bo $A$ $R$-algebra in $V$ $R$-modul.
  \pojem{Upodobitev algebre $A$ nad $V$} je homomorfizem algeber $\rho: A
  \to \End_R(V)$.
\end{definicija}

\begin{opomba}
  Če je $V$ prost $R$-modul, in $V = R^n$, je $\End_R(V) \cong M_n(R)$.
  Številu $n$ pravimo \pojem{stopnja upodobitve}.
\end{opomba}

\begin{definicija}
  Naj bo $R$ komutativen kolobar z enico, $G$ grupa in $V$ $R$-modul.
  \pojem{Upodobitev grupe $G$ nad $V$} je homomorfizem algeber $\rho: RG
  \to \End_R(V)$.
\end{definicija}

